\documentclass[12pt]{article}

% a template that a friend gave, it's worked well enough for me
% i have added some packages and stuff that have proved useful

\usepackage{fancyhdr}
\usepackage{tipa}
\usepackage{fontspec}
\usepackage{amsfonts}
\usepackage{enumitem}
\usepackage[margin=1in]{geometry}
\usepackage{graphicx}
\usepackage{float}
\usepackage{amsmath}
\usepackage{braket}
\usepackage{amssymb}
\usepackage{booktabs}
\usepackage{hyperref}
\usepackage{mathtools}
\usepackage{xcolor}
\usepackage{float}
\usepackage{algpseudocodex}
\usepackage{titlesec}
\usepackage{bbm}

\pagestyle{fancy}
\fancyhf{} % sets both header and footer to nothing
\lhead{Kevin Sheng}
\setmainfont{Comic Neue}
\renewcommand{\headrulewidth}{1pt}
\setlength{\headheight}{0.75in}
\setlength{\oddsidemargin}{0in}
\setlength{\evensidemargin}{0in}
\setlength{\voffset}{-.5in}
\setlength{\headsep}{10pt}
\setlength{\textwidth}{6.5in}
\setlength{\headwidth}{6.5in}
\setlength{\textheight}{8in}
\renewcommand{\headrulewidth}{0.5pt}
\renewcommand{\footrulewidth}{0.3pt}
\setlength{\textwidth}{6.5in}
\usepackage{setspace}
\usepackage{multicol}
\usepackage{float}
\setlength{\columnsep}{1cm}
\setlength\parindent{24pt}
\usepackage [english]{babel}
\usepackage [autostyle, english = american]{csquotes}
\MakeOuterQuote{"}

\setlength{\parskip}{6pt}
\setlength{\parindent}{0pt}

\titlespacing\section{0pt}{12pt plus 4pt minus 2pt}{0pt plus 2pt minus 2pt}
\titlespacing\subsection{0pt}{12pt plus 4pt minus 2pt}{0pt plus 2pt minus 2pt}
\titlespacing\subsubsection{0pt}{12pt plus 4pt minus 2pt}{0pt plus 2pt minus 2pt}

\hypersetup{colorlinks=true, urlcolor=blue}

\newcommand{\correction}[1]{\textcolor{red}{#1}}

\usepackage{nicematrix}

\rhead{Math 110A}

\makeatletter
\def\@seccntformat#1{%
  \expandafter\ifx\csname c@#1\endcsname\c@section\else
  \csname the#1\endcsname\quad
  \fi}
\makeatother

\newcommand{\lra}{\xLeftrightarrow}
\newcommand{\ra}{\xRightarrow}
\DeclareMathOperator{\pr}{pr}
\DeclareMathOperator{\Mat}{Mat}

\begin{document}

\section{Problem 1}

\subsection{Part A}\label{sec:p1a}

Here's the pasted proof of the same thing from last homework:

{
\color{blue}
If we have an $\underline{f}$, then $\underline{f}(\pi(x))=f(x)$ for all $x \in R$.

Consider two functions $f_1$ and $f_2$ that both satisfy this property.
For any $x \in R$, we have that $f_1(\pi(x))=f(x)=f_2(\pi(x))$.
Since $\pi$ is surjective, $f_1(s)=f_2(s)$ for all $s \in R/J$.
}

\subsection{Part B}\label{sec:p1b}

The proof is basically the same as the proof in 3b in the last homework,
but I'll just put it here again for posterity.

I'll prove the backward direction and its converse.

If $f$ annihilates $J$, then $j \in J \implies f(j)=0$.

Let us define $\underline{f}(\pi(v))=f(v)$.
This map is well-defined since $\pi$ is a surjection onto $R/J$ and
\begin{align*}
    \pi(v_1)=\pi(v_2)
        & \implies v_1-v_2 \in J   \\
        & \implies f(v_1-v_2)=0    \\
        & \implies f(v_1)-f(v_2)=0 \\
        & \implies f(v_1)=f(v_2)
\end{align*}
By construction, this satisfies the commutativity diagram as well.

For the backwards direction, BWOC say $\exists j \in J: f(j) \ne 0$.

Since $J$ is a proper two-sided ideal, $R \setminus J$ must be nonempty.
Take some $r \in R \setminus J$.
By definition, $\pi(r)=\pi(r + j)$, so $\underline{f}(\pi(r))=\underline{f}(\pi(r+j))$.
However, $f(r+j)=f(r)+f(j) \ne f(r)$, since the two
being equal would imply $f(s)=0$.
This forces $\underline{f}(\pi(r)) \ne \underline{f}(\pi(r+j))$,
which contradicts the nature of $\underline{f}$. $\square$

\subsection{Part C}

Notice that $\pi^*$ simply takes a given $\underline{f}$ and returns the associated $f$.

From \ref{sec:p1a}, we know that $\pi^*$ is injective.

As for the second part, for any $f$ that annihilates $J$,
we can take the associated $\underline{f}$ that we know must exist by \ref{sec:p1b}
and plug it into $\pi^*$ to get back our $f$.
\ref{sec:p1b} also tells us that if an $f$ is in the range of $\pi^*$, a ring homomorphism that
when plugged into $\pi^*$ gives $f$ can only exist when $f$ annihilates $J$.
Thus, the range of $\pi^*$ is exactly all $f: R \to S$ that annihilate $J$. $\square$

\section{Problem 2}

Let's first show $(r) \subseteq (u \cdot r)$.
Fix any $a, b \in R$, and consider $arb \in (r)$.
Since $u$ is a unit, $u^{-1}$ exists.
If we take the left multiple as $au^{-1}$ and right multiple as $b$,
we see that $arb=\left(au^{-1}\right)(u \cdot r)(b) \in (u \cdot r)$.

The other direction is easier.
Again, fix any $a, b \in R$ and consider $a(u \cdot r)b \in (u \cdot r)$.
Notice that we can reexpress this as $(au)r(b)$ to see that the element is in $(r)$ as well.

The proof for $(r)=(r \cdot u)$ is basically the same.
I won't bore you with that. $\square$

\section{Problem 3}\label{sec:p3}

If $R$ is a field, then every nonzero element is a unit and has a multiplicative inverse.
Consider any two-sided ideal $J$ with a nonzero element $r$.
By closure properties, $r \cdot r^{-1} = 1 \in J$.
Now, for any $x \in R$ we have $x \cdot 1 = x \in J \implies R \subseteq J$ and $R=J$.
Given that any two-sided ideal with a nonzero element must be equivalent to $R$
itself, we can deduce that the only two-sided ideals of $R$ are $\{0\}$ and $R$ itself.

Now consider the case where $R$ wasn't a field.
In other words, there's a nonzero $r \in R$ that doesn't give you $1$
no matter what you multiply it with.

Notice that $(r)=\{xr \mid x \in R\} \notin \{R, \{0\}\}$.
$(r) \ne R$, since $r$ isn't a unit and thus $1 \notin R$.
$(r) \ne \{0\}$ either, since $r \in (r)$.
Thus, there are more two-sided ideals than just $R$ and $\{0\}$. $\square$

\pagebreak

\section{Problem 4}

\subsection{Part A}

Consider any two-sided ideal $J \subseteq R$ that contains a nonzero matrix $M$.

Take $i, j$ s.t. $M_{ij} = x \ne 0$.
Now consider the following matrices that must also be in $J$ by the closure properties:
\begin{enumerate}
    \item If we multiply $M$ by some permutation matrices, we see that
            there must exist some $M' \in J$ s.t. $M'_{11}=x$.

    \item Let $B=AM'A$, where
            \[A_{ij}=\begin{cases}
                    1 & i=j=1            \\
                    0 & \text{otherwise}
                \end{cases}\]
            to get the following matrix:
            \[B_{ij}=(AM'A)_{ij}=\begin{cases}
                    x & i=j=1            \\
                    0 & \text{otherwise}
                \end{cases}\]
    
    \item Since $B \in J$, $O=B \cdot \left(x^{-1} \cdot I_{n \times n}\right) \in J$ as well,
    where $I$ is the identity matrix.
    Since the matrix we're multiplying $B$ by scales every element by $x^{-1}$,
    this gives a matrix with a $1$ in the top right corner and nothing else.

    \item Multiplying $O$ by some permutation matrices means that any $O'_p$ defined by
    \[(O'_p)_{ij}=\begin{cases}
        1 & i=j=p \\
        0 & \text{otherwise}
    \end{cases}\]
    must be in $J$ as well.
    
    \item Adding up $O'_1$ through $O'_n$ gives that $I_{n \times n} \in J$.
\end{enumerate}
Since we've proven that the identity matrix is now in $J$, our two-sided
ideal must be equal to $R$ by similar reasoning as in \ref{sec:p3}.

\pagebreak

\subsection{Part B}

If $k$ isn't a field, take a nonzero element $c \in k$ that doesn't have an inverse.

Consider $\left(c \cdot I_{n \times n}\right)$.
It's not equal to $\{0\}$, since it has $cI$.
To show that $(cI) \ne R$, I'll show that $I \notin (cI)$.

BWOC say $I \in cI$.
This means we can express it like so:
\[I=\sum_{i=1}^{m} A_i(cI)B_i=c\left(\sum_{i=1}^{m} A_i \cdot I \cdot B_i\right)\]
This is a contradiction, multiplying any matrix by $c$ can't result in $I$.
This follows from that $c$ doesn't have a multiplicative inverse.

\subsection{Part C}

The two previous parts of this problem extended the results of part 3
to a certain case of non-commutative rings.

\section{Problem 5}

\begin{enumerate}
    \item $n=2$: Yes. $1 \cdot 1 \equiv 1\pmod 2$.

    \item $n=3$: Yes. $2 \cdot 2\equiv 1\pmod 3$, and $1$ is still its own inverse.

    \item $n=4$: No. $2 \cdot 1 \equiv 2$, $2 \cdot 2 \equiv 0$, and $2 \cdot 3\equiv 2\pmod 4$.
            No element gives $1$ when you multiply it by $2$.

    \item $n=5$: Yes. $2 \cdot 3\equiv 1$, $3 \cdot 2 \equiv 1$, and $4 \cdot 4 \equiv 1$.

    \item $n=6$: No. $3 \cdot 2 \equiv 0$, $3 \cdot 3 \equiv 3$, $3 \cdot 4 \equiv 0$, and $3 \cdot 5 \equiv 5$.
\end{enumerate}

\section{Problem 6}

As proven in class, the four elements of $\Z/4\Z$ are $0$, $1$, $2$, and $3$.

The elements of $\Z/2\Z[x]$ are all polynomials with coefficients $0$ or $1$.
If we make $x^2 \equiv 0$, we essentially limit ourselves to all polynomials
of degree $1$ or lower, limiting ourselves to these four elements:
\[x+1, x, 1, 0\]

Any hypothetical isomorphism of these two rings would have $0$ and $1$ map to each other.

However, if we let $\phi: \Z/4\Z \to (\Z/2\Z)[x]/(x^2)$, then $\phi(1+1)=\phi(1)+\phi(1)=0$,
which immediately bans any sort of isomorphism as the forced homomorphism is not injective.

\end{document}
